\documentclass[10pt, letterpaper]{article}

% Pacotes e configuração visual
\usepackage[
    ignoreheadfoot,
    top=1.1 cm,
    bottom=1.1 cm,
    left=1.3 cm,
    right=1.3 cm,
    footskip=0.6 cm,
]{geometry}
\usepackage{titlesec}
\usepackage{tabularx}
\usepackage{array}
\usepackage[dvipsnames]{xcolor}
\definecolor{primaryColor}{RGB}{0, 0, 0}
\usepackage{enumitem}
\usepackage{fontawesome5}
\usepackage{amsmath}
\usepackage[
    pdftitle={Vitor Jacom - Currículo},
    pdfauthor={V\'itor Jacom},
    colorlinks=true,
    urlcolor=primaryColor
]{hyperref}
\usepackage[pscoord]{eso-pic}
\usepackage{calc}
\usepackage{bookmark}
\usepackage{lastpage}
\usepackage{changepage}
\usepackage{paracol}
\usepackage{ifthen}
\usepackage{needspace}
\usepackage{iftex}

\ifPDFTeX
    \input{glyphtounicode}
    \pdfgentounicode=1
    \usepackage[T1]{fontenc}
    \usepackage[utf8]{inputenc}
    \usepackage{lmodern}
\fi

\usepackage{charter}

\raggedright
\AtBeginEnvironment{adjustwidth}{\partopsep0pt}
\pagestyle{empty}
\setcounter{secnumdepth}{0}
\setlength{\parindent}{0pt}
\setlength{\topskip}{0pt}
\setlength{\columnsep}{0.15cm}
\pagenumbering{gobble}

\titleformat{\section}{\needspace{4\baselineskip}\bfseries\large}{}{0pt}{}[\vspace{1pt}\titlerule]
\titlespacing{\section}{-1pt}{0.3 cm}{0.2 cm}
\renewcommand\labelitemi{$\vcenter{\hbox{\small$\bullet$}}$}

\newenvironment{highlights}{
    \begin{itemize}[
        topsep=0.10 cm,
        parsep=0.10 cm,
        partopsep=0pt,
        itemsep=0pt,
        leftmargin=0 cm + 10pt
    ]
}{
    \end{itemize}
}

\newenvironment{header}{
    \setlength{\topsep}{0pt}\par\kern\topsep\centering\linespread{1.5}
}{
    \par\kern\topsep
}

\begin{document}

% ====== Cabeçalho ======
\begin{header}
    {\fontsize{25 pt}{25 pt}\selectfont V\'itor Jacom}

    \vspace{5 pt}

    \normalsize
    Porto Alegre -- RS -- Brasil \quad \href{mailto:vjds.vitor@gmail.com}{vjds.vitor@gmail.com} \quad \href{tel:+55 51 99578 3576}{(51) 99578-3576} \\
    \href{https://www.linkedin.com/in/vitorjacom/}{linkedin.com/in/vitorjacom} \quad \href{https://github.com/VitorJacom}{github.com/VitorJacom}
\end{header}

% ====== Resumo Profissional ======
\section*{Resumo Profissional}
\textbf{Engenheiro de Software Full-Stack} com experi\^encia pr\'atica nos ecossistemas \textbf{Java (Spring Boot)} e \textbf{TypeScript (React, React Native, NestJS)}. Atua\c{c}\~ao do desenvolvimento de \textbf{APIs RESTful} \'a implanta\c{c}\~ao em nuvem (\textbf{AWS e GCP}). Dom\'inio de \textbf{Docker}, esteiras \textbf{CI/CD}, automa\c{c}\~ao de testes e arquitetura (\textbf{DDD}, Clean Code), focado em alta performance, qualidade e entregas de valor de neg\'ocio.

% ====== Tecnologias ======
\section*{Compet\^encias T\'ecnicas}
\noindent\textbf{Linguagens:} Java (11/17+), TypeScript, JavaScript, Kotlin, Python, SQL \\
\textbf{Front-end \& Mobile:} React.js, React Native, Next.js, HTML5, CSS3, Tailwind CSS \\
\textbf{Back-end \& Frameworks:} Spring Boot, Spring Security, NestJS, Node.js, Express, APIs RESTful \\
\textbf{Cloud \& DevOps:} AWS (Lambda, Aurora, S3, ECR), Google Cloud (Cloud Run), Docker, GitHub Actions, GitLab CI/CD \\
\textbf{Bancos de Dados \& ORM:} PostgreSQL, MySQL, Redis, Prisma ORM, Spring Data JPA \\
\textbf{Testes \& Qualidade:} JUnit, Mockito, Jest, Cypress, Maestro, Postman

% ====== Experiência ======
\section*{Experi\^encia Profissional}

\textbf{Desenvolvedor (Freelancer)} \hfill Remoto \hfill (Nov 2025 -- Atual)

\begin{itemize}
    \item \textbf{Desenvolvimento Full-Stack, TS/React \& CMS:} Projetou e desenvolveu do zero portais institucionais utilizando TypeScript e React, incluindo o desenvolvimento de um painel CMS interno para gest\~ao aut\^onoma de conte\'udo (membros da equipe, artigos e not\'icias), formul\'ario de contato integrado por e-mail e modais de aviso antifraude.

    \item \textbf{Design Mobile-First \& Responsividade:} Projetou interfaces altamente responsivas focadas em performance e navega\c{c}\~ao mobile-first, garantindo usabilidade impec\'avel para usu\'arios em dispositivos m\'oveis (perfil predominante dos clientes finais) e tempos de carregamento reduzidos.

    \item \textbf{Integra\c{c}\~oes de APIs \& Redes Sociais:} Implementou integra\c{c}\~ao com a Meta Graph API para sincroniza\c{c}\~ao e exibi\c{c}\~ao autom\'atica no site das \'ultimas publica\c{c}\~oes do Instagram, mantendo o portal din\^amico sem necessidade de interven\c{c}\~ao manual.

    \item \textbf{Migra\c{c}\~ao de Dom\'inios, Infraestrutura \& SEO:} Realizou a migra\c{c}\~ao de infraestruturas legadas (Wix) para solu\c{c}\~oes customizadas, gerenciando a transfer\^encia de dom\'inios (Registro.br), zonas DNS e indexa\c{c}\~ao SEO via Google Search Console para alavancar a visibilidade org\^anica.

    \item \textbf{Engenharia de Requisitos \& Gest\~ao de Projetos:} Conduziu autonomamente todo o ciclo de vida do software (SDLC), desde a idea\c{c}\~ao de mockups e levantamento de requisitos com clientes at\'e o planejamento de entregas semanais e alinhamento direto com os stakeholders.
\end{itemize}

\textbf{Software Engineer Intern} \hfill Tecnova Energia -- Porto Alegre, RS \hfill (Abr 2024 -- Out 2025)

\begin{itemize}
    \item \textbf{Engenharia de Dados \& Cloud Serverless:} Eliminou a migra\c{c}\~ao manual de dados de planejamento (+25k linhas) para o Power BI, construindo uma aplica\c{c}\~ao serverless conteinerizada com Docker e implantada no GCP (Google Cloud Run) sob arquitetura DDD, garantindo integridade dos dados e custo zero em per\'iodos de ociosidade.

    \item \textbf{Automa\c{c}\~ao ERP \& CTM:} Desenvolveu microsservi\c{c}os em Java 17 (Spring Boot) para integra\c{c}\~ao cont\'inua entre o sistema de Gest\~ao de Viagens Corporativas (CTM -- Corporate Travel Management) e o ERP (Enterprise Resource Planning / Gest\~ao Empresarial), automatizando a aprova\c{c}\~ao de despesas e a sincroniza\c{c}\~ao or\c{c}ament\'aria em tempo real.

    \item \textbf{Pipeline de ETL \& Engenharia de Dados:} Projetou e implementou pipelines de ingest\~ao para extrair, normalizar e estruturar dados originados de ferramentas de engenharia como greendocs. A automa\c{c}\~ao padronizou dados n\~ao estruturados no banco relacional, alimentando dashboards executivos em BI.

    \item \textbf{Pesquisa T\'ecnica \& DevOps:} Estruturou rotinas de CI/CD via GitHub Actions e padronizou ambientes locais com Docker, atuando com autonomia na interpreta\c{c}\~ao de documenta\c{c}\~oes oficiais para construir integra\c{c}\~oes e APIs RESTful em Java e Python.
\end{itemize}

% ====== Projetos e Consultoria ======
\section*{Portf\'olio de Engenharia de Software (F\'abrica de Software -- AGES/PUCRS)}
\textit{\small Projetos desenvolvidos em contexto acad\^emico-profissional para clientes reais, utilizando metodologias \'ageis e squads multidisciplinares.}
\vspace{0.3em}

\textbf{Estudify (Cliente: Startup)} \hfill (2026) \\
Aplicativo educacional gamificado focado em estudos para o ENEM e vestibulares.
\begin{itemize}
    \item \textbf{Governan\c{c}a \& Gest\~ao de Squads (AGES IV):} Liderou e coordenou a governan\c{c}a t\'ecnica de uma equipe de 18 pessoas distribu\'ida uniformemente em 4 squads simult\^aneas, estruturando refinamentos de User Stories, prioriza\c{c}\~ao de backlog no ClickUp e entregas quinzenais alinhadas diretamente com o cliente startup.

    \item \textbf{Engenharia de Qualidade em Hardware Real (V\&V):} Projetou a estrat\'egia de automa\c{c}\~ao de testes E2E mobile integrando Maestro e ADB (Android Debug Bridge) executados em um dispositivo f\'isico Android real, e criou um servidor Node.js auxiliar para limpeza idempotente do banco PostgreSQL, eliminando inconsist\^encias de ambiente de testes.

    \item \textbf{Otimiza\c{c}\~ao de Performance \& Builds:} Otimizou o gerenciamento de ativos gr\'aficos e compila\c{c}\~ao do projeto, reduzindo em 98\% o consumo de mem\'oria no Android Studio, sanando erros de estouro de mem\'oria no compilador e garantindo a gera\c{c}\~ao est\'avel de APKs de produ\c{c}\~ao.

    \item \textbf{Stack T\'ecnica:} \textbf{React Native}, \textbf{TypeScript}, \textbf{NestJS}, \textbf{Prisma ORM}, \textbf{PostgreSQL}, \textbf{Docker}, \textbf{Maestro}, \textbf{ADB}.
\end{itemize}

\textbf{Lucky Draw (Cliente: Startup)} \hfill (2025) \\
Aplicativo mobile para desbloqueio criativo e gera\c{c}\~ao de desafios di\'arios de desenho.
\begin{itemize}
    \item \textbf{Arquitetura Cloud-Native Serverless:} Projetou e implementou uma arquitetura serverless na AWS utilizando AWS Lambda (backend Spring Boot), AWS Aurora PostgreSQL para dados relacionais e AWS S3 para armazenamento de m\'idias/desenhos dos usu\'arios, zerando custos em per\'iodos de ociosidade.

    \item \textbf{DevOps \& Pipeline CI/CD Automatizada:} Configurou esteira de CI/CD com GitLab CI cobrindo est\'agios de build, execu\c{c}\~ao de testes (JUnit/Mockito), an\'alise est\'atica de c\'odigo com Checkstyle e empacotamento em cont\^eineres Docker enviados para o AWS ECR.

    \item \textbf{Governan\c{c}a de C\'odigo \& Code Review:} Atuou como revisor t\'ecnico principal de Merge Requests (MRs) da squad, garantindo a ader\^encia ao padr\~ao Controller-Service-Repository, conven\c{c}\~oes RESTful, cobertura de testes e fluxo do GitFlow.

    \item \textbf{Seguran\c{c}a, Tratamento de Exce\c{c}\~oes \& APIs:} Desenvolveu rotas RESTful protegidas por Spring Security + JWT, configurou manipulador global de exce\c{c}\~oes (GlobalExceptionHandler) e padronizou respostas e valida\c{c}\~oes com Bean Validation (@Valid) e Swagger.

    \item \textbf{Stack T\'ecnica:} \textbf{Java 17}, \textbf{Spring Boot}, \textbf{PostgreSQL}, \textbf{AWS (Lambda, Aurora, S3, ECR)}, \textbf{Docker}, \textbf{GitLab CI/CD}, \textbf{JUnit}, \textbf{Mockito}, \textbf{Checkstyle}.
\end{itemize}

\textbf{MeditAmamente (Cliente: Pesquisa Acad\^emica para uma Doutoranda)} \hfill (2024) \\
Plataforma digital para suporte \`a sa\'ude mental e protocolo de medita\c{c}\~ao guiada para m\~aes de beb\^es prematuros.
\begin{itemize}
    \item \textbf{Arquitetura Clean \& DDD (Lideran\c{c}a T\'ecnica):} Atuando al\'em do n\'ivel formal, projetou e implementou a arquitetura back-end do sistema em Spring Boot, estruturando uma abordagem em 4 camadas baseada em Domain-Driven Design (DDD) (Domain, Application, Infrastructure e Presentation) para garantir desacoplamento total e facilidade de manuten\c{c}\~ao.

    \item \textbf{Seguran\c{c}a, Autentica\c{c}\~ao \& Configura\c{c}\~oes:} Implementou o m\'odulo de seguran\c{c}a com Spring Security e tokens JWT para autentica\c{c}\~ao e autoriza\c{c}\~ao stateless, al\'em de estruturar o mapeamento de entidades relacionais e configura\c{c}\~oes no PostgreSQL.

    \item \textbf{Garantia de Qualidade \& Refatora\c{c}\~ao:} Conduziu refatora\c{c}\~oes para elimina\c{c}\~ao de d\'ebitos t\'ecnicos e padroniza\c{c}\~ao de tratamento de erros no back-end, colaborando diretamente no alinhamento arquitetural do time.

    \item \textbf{Reconhecimento \& Premia\c{c}\~ao:} O projeto obteve a nota m\'axima e foi premiado como Projeto Destaque Acad\^emico PUCRS, superando crit\'erios rigorosos de arquitetura, entrega t\'ecnica e valor gerado para a pesquisa.

    \item \textbf{Stack T\'ecnica:} \textbf{Java}, \textbf{Spring Boot}, \textbf{PostgreSQL}, \textbf{Domain-Driven Design (DDD)}, \textbf{Clean Architecture}, \textbf{Spring Security}, \textbf{JWT}, \textbf{Docker}, \textbf{Git}.
\end{itemize}

\textbf{Globo Aplausos (Cliente: Globo\texttrademark)} \hfill (2023) \\
Plataforma interna de gamifica\c{c}\~ao e reconhecimento corporativo desenvolvida em parceria com a Globo.
\begin{itemize}
    \item \textbf{Desenvolvimento Full-Stack (AGES I):} Desenvolveu componentes responsivos de interface com Next.js, React, TypeScript e Material UI (MUI), integrando com endpoints RESTful em NestJS e realizando o mapeamento de entidades via Prisma ORM com banco MySQL.

    \item \textbf{Automa\c{c}\~ao no Banco de Dados:} Projetou o motor de gamifica\c{c}\~ao do sistema utilizando MySQL Event Triggers para automatizar a distribui\c{c}\~ao peri\'odica de moedas virtuais aos funcion\'arios, validando e isolando o ambiente de banco de dados via Docker.

    \item \textbf{Qualidade de C\'odigo \& Testes Unit\'arios:} Escreveu su\'ites de testes unit\'arios para servi\c{c}os e controllers do back-end utilizando Jest, assegurando a integridade das regras de neg\'ocio, al\'em de integrar testes End-to-End (E2E) com Cypress para fluxos cr\'iticos.

    \item \textbf{Padroniza\c{c}\~ao \& Governan\c{c}a \'Agil:} Participou de uma squad \'agil.

    \item \textbf{Stack T\'ecnica:} \textbf{Next.js}, \textbf{React}, \textbf{TypeScript}, \textbf{NestJS}, \textbf{Prisma ORM}, \textbf{MySQL}, \textbf{Material UI (MUI)}, \textbf{Jest}, \textbf{Cypress}, \textbf{Husky}, \textbf{Docker}, \textbf{Git}.
\end{itemize}

% ====== Formação ======
\section*{Forma\c{c}\~ao Acad\^emica}
\textbf{Bacharel em Engenharia de Software} -- PUCRS \hfill (Conclu\'ido em Ago/2026)
\begin{highlights}
    \item \textbf{Projeto Destaque Acad\^emico (2024):} Premia\c{c}\~ao m\'axima pelo desenvolvimento da plataforma MeditAmamente.
    \item \textbf{Monitoria de POO em Java (2023):} Aux\'ilio a alunos em l\'ogica de programa\c{c}\~ao e orienta\c{c}\~ao a objetos.
    \item \textbf{Instrutor T\'ecnico (2025):} Ministrou o workshop "Containers e Docker" na Semana Acad\^emica da PUCRS.
\end{highlights}

% ====== Certificações e Idiomas ======
\section*{Certifica\c{c}\~oes e Idiomas}

\begin{itemize}[leftmargin=0.35cm,itemsep=1pt,parsep=0pt,topsep=2pt]
    \item \textbf{GitHub Foundations (GH-900)} -- GitHub (2026)
    \item \textbf{Ingl\^es:} N\'ivel Avan\c{c}ado
\end{itemize}

\end{document}